\chapter{REVISIÓN DE LITERATURA}
\label{ch:revision}

Este capítulo revisa la literatura relevante para fundamentar la propuesta metodológica de la memoria. Se organiza en dos secciones: primero, la evidencia empírica sobre eficiencia
y desempeño en economía social y cooperativas financieras; segundo, las
experiencias internacionales en construcción de cuentas satélite para la
economía social, incluyendo la situación del contexto latinoamericano y
chileno.

\section{Economía social y organizaciones de la sociedad civil}

La Economía Social y Solidaria (ESS) agrupa a organizaciones cuyo funcionamiento 
se basa en principios de solidaridad, autogestión democrática y primacía del ser 
humano sobre el capital ---concepto que se desarrolla en detalle en el Marco 
Teórico---. Contar con marcos institucionales claros y sistemas de información 
robustos resulta fundamental para evaluar el desempeño de estas organizaciones.
Un aspecto especialmente relevante es la \textbf{eficiencia de las organizaciones no
gubernamentales (ONGs)}, para la cual se han desarrollado metodologías orientadas a
analizar el uso de recursos y el retorno social generado \cite{coupet_toward_2019}.
Estos enfoques permiten comparar el aporte de las organizaciones de la sociedad civil
(OSC) en distintos contextos y constituyen el punto de partida para comprender cómo
medir el desempeño de organizaciones con objetivos que van más allá del lucro.

En el ámbito de las \textbf{cooperativas financieras}, la eficiencia depende no solo
del desempeño económico, sino también de la capacidad de generar valor social y
sostener estructuras de gobernanza democráticas \shortcite{perilleux_are_2016}. Este
equilibrio entre objetivos económicos y sociales constituye una característica
distintiva de las organizaciones de la economía social y plantea desafíos particulares
para su medición dentro de los marcos estadísticos convencionales, aspecto que se
desarrolla en la subsección siguiente. Para el caso chileno, \cite{lemus_credit_2022}
proveen evidencia econométrica de que las cooperativas de ahorro y crédito cumplen un
rol de inclusión financiera, atendiendo en mayor proporción a personas de bajos
ingresos, adultos mayores, mujeres y habitantes de localidades alejadas de los
grandes centros urbanos.

\subsection{Cooperativas financieras y medición de valor agregado}

Las \textbf{cooperativas de ahorro y crédito (CAC)} son organizaciones financieras
de propiedad y control democrático de sus socios, cuyo objeto central es la
intermediación financiera al servicio de sus propios miembros. Su medición económica
presenta desafíos específicos dentro de los marcos estadísticos convencionales, ya
que a diferencia de los bancos comerciales operan bajo una lógica de excedentes
limitados que dificulta la aplicación directa de los mismos criterios de estimación
\shortcite{perilleux_are_2016}.

La medición del impacto económico de las cooperativas ha sido abordada en la
literatura a través de distintos enfoques. \cite{uzea_challenges_2015} identifican
tres desafíos principales: la heterogeneidad de las formas jurídicas cooperativas,
la ausencia de registros estandarizados que permitan comparaciones entre
organizaciones, y la dificultad de separar la producción de mercado de la producción
no de mercado en organizaciones que sirven simultáneamente a sus socios y a la
comunidad. Estos desafíos son especialmente relevantes en el contexto chileno, donde
las CAC operan bajo marcos regulatorios distintos según su tamaño patrimonial.

Para abordar estas limitaciones, \shortcite{brown_co-operatives_2015} proponen un marco
de medición del impacto cooperativo que combina indicadores financieros ---como el
valor agregado bruto, el empleo y las remuneraciones--- con indicadores sociales que
capturan dimensiones no monetarias del desempeño organizacional. Este enfoque es
coherente con los lineamientos del manual de Naciones Unidas
\cite{vereinte_nationen_satellite_2018}, que reconoce la naturaleza dual de las
cooperativas como unidades productivas y como organizaciones con misión social.

La literatura especializada distingue además entre organizaciones cuya producción se
valora al costo ---como las organizaciones sin fines de lucro--- y aquellas que
generan producción de mercado directamente incorporable al Sistema de Cuentas
Nacionales \shortcite{perilleux_are_2016,barea_tejeiro_manual_2007}. Esta distinción tiene
consecuencias directas sobre el procedimiento de estimación aplicable a cada tipo de
organización, aspecto que se desarrolla en el Marco Teórico.


\section{Experiencias internacionales y sistemas de información para organizaciones
sociales}

Esta sección revisa la evidencia internacional sobre la medición del aporte económico
de la economía social. Se organiza en tres partes: primero se discute el problema de
la invisibilidad estadística y cómo las cuentas satélite lo abordan; luego se revisan
las experiencias de España, Portugal y Polonia como casos de referencia; y finalmente
se analiza el contexto latinoamericano y chileno, incluyendo los desafíos de
información y los requerimientos metodológicos que se derivan de la experiencia
comparada.

\subsection{Invisibilidad estadística y cuentas satélite}

Las organizaciones de la economía social han sido históricamente subrepresentadas en
las estadísticas oficiales, ya que sus actividades se diluyen dentro de los marcos
tradicionales de cuentas nacionales \cite{barea_tejeiro_manual_2007}. Esta limitación
dificulta la medición de su contribución económica y social y ha motivado el desarrollo
de herramientas estadísticas específicas para el sector, entre las que destacan las
cuentas satélite ---cuyo funcionamiento se desarrolla en el Marco Teórico---. La
experiencia comparada, que se revisa a continuación, muestra que esta metodología ha
sido implementada con éxito en distintos contextos institucionales, aunque siempre
enfrentando desafíos estructurales de información similares.

\subsection{Experiencias internacionales}

La experiencia más reciente y metodológicamente detallada proviene de España. En 2024,
el Instituto Nacional de Estadística publicó la Cuenta Satélite de la Economía Social
cubriendo el período 2019--2023 \cite{noauthor_cuenta_2024}. El proceso articuló
cuatro fuentes: el Directorio Central de Empresas (DIRCE) para identificar entidades,
los registros de la Tesorería General de la Seguridad Social para empleo y
remuneraciones, las declaraciones del Impuesto sobre Sociedades para variables
económicas de producción y consumo intermedio, y encuestas del Ministerio de Trabajo
para el voluntariado. El identificador fiscal (NIF) operó como variable de enlace
entre todas las bases, equivalente al RUT en el contexto chileno. La distinción clave
del caso español es que la declaración tributaria entrega valores exactos de ingresos
y gastos por empresa, lo que permite calcular producción y consumo intermedio
directamente sin recurrir a proxies \cite{noauthor_cuenta_2024}.

En una línea similar, aunque con algunas diferencias metodológicas relevantes,
Portugal desarrolló su \textit{Conta Satélite da Economia Social} (CSES) en
colaboración entre el Instituto Nacional de Estatística y la Cooperativa António
Sérgio para la Economía Social (CASES), publicando su cuarta edición en 2023 con
datos para 2019--2020 \shortcite{pedroso_conta_2023}. Portugal fue pionero en adoptar el
Manual ONU TSE de 2018 como referencia, reclasificando sus series históricas para
ganar comparabilidad internacional. Sus resultados muestran que la economía social
representó el 3{,}2\,\% del VAB nacional en 2020, con 73.851 entidades activas
empleando el 5{,}0\,\% del empleo total. Para las variables económicas, Portugal
empleó el sistema de Información Empresarial Simplificada (IES) del Ministerio de
Justicia ---que consolida estados financieros anuales--- junto con el Relatório Único
del Ministerio de Trabajo para remuneraciones, y estados financieros de cooperativas
credenciadas por CASES para el sector cooperativo específicamente.

Un tercer caso de referencia es el de Polonia, que publicó en 2021 su primera cuenta
satélite con datos de 2018, ejecutada por Statistics Poland en el marco de un proyecto
Eurostat \cite{statistics_poland_social_2021}. El caso polaco es especialmente
relevante porque su entorno estadístico era más limitado: la encuesta anual de
empresas (SP) no cubría cooperativas pequeñas ni sus funciones sociales, lo que obligó
a desarrollar una encuesta especial (ES-S) aplicada por primera vez en 2017. Esta
encuesta se complementó con datos del impuesto corporativo (CIT-8) y del registro de
seguridad social (ZUS) para completar las variables de producción, consumo intermedio
y remuneraciones. Para el voluntariado, Polonia implementó una encuesta de trabajo no
remunerado (PNZ) como módulo adicional de su Encuesta de Fuerza Laboral.

A pesar de sus diferencias de contexto, los tres países comparten el mismo patrón
estructural: un registro de entidades cruzado con fuentes económicas mediante
identificador fiscal, con encuestas complementarias para lo que los registros
administrativos no capturan. La experiencia comparada revela además que ningún país
ha cubierto todos los bloques de información requeridos con fuentes preexistentes en
una primera implementación. La estrategia estándar ha sido construir una versión
inicial con los datos disponibles, documentar las brechas explícitamente, y proponer
una agenda de mejora para versiones posteriores \cite{archambault_lester_2022}. Esta
lógica iterativa es la que orienta la propuesta metodológica de esta memoria.

\subsection{Contexto latinoamericano y chileno}

A diferencia de la trayectoria europea, en América Latina la implementación de cuentas
satélite para la economía social se encuentra en etapas iniciales. México constituye
el caso más avanzado de la región: el INEGI publicó un estudio de caso sobre la
economía social mexicana para los años 2013 y 2018, utilizando como referencia los
lineamientos del SCN 2008 y el Manual CIRIEC, con el apoyo financiero del FIDA a
través de la CEPAL \cite{ciriec-mexico_estudio_2022}. Los resultados indican que en
2018 las sociedades cooperativas de ahorro y préstamo representaron el 15{,}1\,\% del
PIB de la economía social mexicana. Metodológicamente, el estudio empleó el Registro
Público de Comercio para delimitar la población y encuestas sectoriales para las
variables económicas.

Costa Rica representa un caso de interés por la madurez relativa de su sector
cooperativo y su activa agenda institucional en esta materia. Según datos del IV Censo
Nacional Cooperativo, el cooperativismo costarricense contribuye con el 7\,\% del PIB
nacional, con casi el 40\,\% de la Población Económica Activa vinculada al sector
\cite{oibescoop_costa_2023}. Sin embargo, esta relevancia económica no se ha traducido
en un sistema de medición sistemática: a pesar de contar desde 2019 con un capítulo
nacional del CIRIEC radicado en la Universidad Estatal a Distancia (UNED), una cuenta
satélite de la economía social como tal no ha sido implementada.
\cite{castillo_arias_acercamiento_2022} propone un acercamiento teórico hipotético a
la cuenta satélite para el sector cooperativo costarricense, reconociendo
explícitamente que las cuentas satélite aún no se han implementado para este grupo en
el país. La experiencia costarricense ilustra una brecha que se repite en la región:
sectores cooperativos relevantes y con institucionalidad desarrollada, pero sin los
sistemas estadísticos necesarios para medir su aporte económico de forma comparable.

Esta situación es análoga a la que enfrenta Chile. En el país, la experiencia más
cercana en construcción de cuentas satélite sectoriales es la Cuenta Satélite de
Tecnologías de Información y Comunicación, desarrollada en 2006 por el Ministerio de
Economía en conjunto con el INE \cite{noauthor_cuenta_2006}. Su construcción requirió
una encuesta ad hoc aplicada a empresas del sector, junto con registros de aduanas y
de telecomunicaciones. Esta experiencia demuestra la viabilidad técnica del enfoque en
el contexto nacional, pero también evidencia que las cuentas satélite requieren fuentes
estadísticas de mayor granularidad que las disponibles actualmente para el sector
cooperativo, lo que refuerza la relevancia de desarrollar metodologías adaptadas a
contextos con información administrativa limitada.

